\section*{Bab \ref{ch:g9:atoms-to-galaxies} --- Dari Atom sampai Galaksi: Skala Alam Semesta}

\begin{solution}{exo:g9:atoms-to-galaxies:1}
\omterm{def:g9:atoms-to-galaxies:atom}{Inti} ($10^{-15}$ \unit{m}), \omterm{def:g9:atoms-to-galaxies:atom}{atom} ($10^{-10}$), \omterm{def:g7:states-of-matter:molecule}{molekul} (sekitar
$10^{-9}$), semut ($10^{-2}$, atau $10^{-3}$ untuk yang sederhana).
\end{solution}

\begin{solution}{exo:g9:atoms-to-galaxies:2}
Sebuah \omterm{def:g9:atoms-to-galaxies:atom}{inti} yang sangat kecil, \omterm{def:g4:weight-and-mass:weight}{berat}, dan bermuatan positif di
pusatnya; sebuah awan \omterm{def:g9:atoms-to-galaxies:atom}{elektron} negatif yang hampir tanpa bobot jauh di
luarnya. Ukurannya: \omterm{def:g9:atoms-to-galaxies:atom}{atomnya} dekat $10^{-10}$ \unit{m}, \omterm{def:g9:atoms-to-galaxies:atom}{intinya} dekat
$10^{-15}$ --- selisih seratus ribu kali lipat yang membuat
kacang-polong-di-dalam-stadionnya.
\end{solution}

\begin{solution}{exo:g9:atoms-to-galaxies:3}
\omterm{def:g9:atoms-to-galaxies:atom}{Elektron} --- yang negatif, yang berarak dari $-$ ke $+$: berlawanan
dengan perjanjiannya. Hukumnya selamat karena mereka tidak pernah
bergantung pada siapa yang berarak: arusnya, jumlah percabangannya,
potret Ohmnya, dan tagihan \omterm{def:g9:electric-power-energy:power}{dayanya} terbaca serupa dengan anak panah
yang dipertahankan menurut kesepakatan sedunia.
\end{solution}

\begin{solution}{exo:g9:atoms-to-galaxies:4}
Pangkat sepuluh yang terdekat. Tinggi badan: $10^{0}$ \unit{m};
ruang kelas: $10^{1}$; garis tengah Bumi: $10^{7}$.
\end{solution}

\begin{solution}{exo:g9:atoms-to-galaxies:5}
\omterm{def:g9:atoms-to-galaxies:atom}{Atom} ke \omterm{def:g9:atoms-to-galaxies:atom}{inti}: lima anak tangga. Kamu ke Bumi: tujuh. Jarak Matahari
($10^{11}$) ke \omterm{def:g4:solar-system:star}{bintang} terdekat ($4 \times 10^{16}$): lima sampai enam
anak tangga --- beberapa ratus ribu kali lipat.
\end{solution}

\begin{solution}{exo:g9:atoms-to-galaxies:6}
\omterm{def:g9:atoms-to-galaxies:atom}{Atomnya} adalah $10^{-10}$ \unit{m} yang sebagian besar kosong di
sekeliling sebuah \omterm{def:g9:atoms-to-galaxies:atom}{inti} $10^{-15}$: kemasilah \omterm{def:g9:atoms-to-galaxies:atom}{atomnya} menjadi sebuah
meja dan kehampaannya ikut serta. Kepejalannya adalah penolakan
listrik awan \omterm{def:g9:atoms-to-galaxies:atom}{elektronnya} untuk saling bertindih --- kekosongan, yang
dibela dengan kukuh.
\end{solution}

\begin{solution}{exo:g9:atoms-to-galaxies:7}
Milimeternya, $10^{-3}$; \omterm{def:g9:atoms-to-galaxies:atom}{atomnya}, $10^{-10}$: tujuh anak tangga ---
kira-kira sepuluh juta \omterm{def:g9:atoms-to-galaxies:atom}{atom} pada deretnya.
\end{solution}

\begin{solution}{exo:g9:atoms-to-galaxies:8}
\omterm{def:g8:speed-of-light:lightyear}{Tahun cahaya}: $10^{16}$ \unit{m} (sembilan setengah ribu juta juta).
\omterm{ex:g5:stars-night-sky:milkyway}{Bima Sakti}: $10^{21}$. Langit tersurvei yang terdalam: dekat
$10^{26}$.
\end{solution}

\begin{solution}{exo:g9:atoms-to-galaxies:9}
Angka $10^{5}$ kacang polong: $0.01 \times 10^{5} = \qty{1000}{m}$ ---
sebuah ``stadion'' selebar sekilometer: lebih megah daripada stadion
yang sungguhan mana pun, tetapi gambarannya bertahan --- kacang
polongnya di titik pusatnya, dan kacang polong berikutnya sekilometer
jauhnya.
\end{solution}

\begin{solution}{exo:g9:atoms-to-galaxies:10}
Tetes di lautannya: kira-kira $10^{21} \times 10^{5} = 10^{26}$ ---
sehingga pengakuan lamanya, kalau diambil harfiah, berlebihan:
lautannya memuat lebih banyak tetes daripada banyaknya \omterm{def:g7:states-of-matter:molecule}{molekul} yang
dimuat setetes, dengan selisih beberapa anak tangga. Versinya yang
jujur --- \omterm{def:g7:states-of-matter:molecule}{molekul} setetes melebihi jumlah apa pun yang dapat dihitung
seumur hidup --- selamat; tangganya menjaga bahkan pengakuan yang
disayangi pun tetap jujur.
\end{solution}

\begin{solution}{exo:g9:atoms-to-galaxies:11}
Kira-kira sepuluh anak tangga turun sampai ke \omterm{def:g9:atoms-to-galaxies:atom}{atomnya} ($10^{0}$ sampai
$10^{-10}$) dan dua puluh enam naik sampai ke langit yang dalam ---
dekat tengahnya, sedikit condong ke arah yang kecil. Sembilan tahun
persekolahan menempatkan seorang pembaca di dalam jangkauan yang jujur
dari kedua ujung segala yang terukur.
\end{solution}

\begin{solution}{exo:g9:atoms-to-galaxies:12}
Jawabannya bersifat pribadi --- latihannya adalah suratnya sendiri.
(Satu contoh, sekadar untuk bentuknya: kepada pembaca kelas satu pada
bab \omterm{def:g1:light-and-shadows:shadow}{bayang-bayangnya} --- ``si kembar gelap di trotoar suatu hari akan
menjelaskan gerhana Matahari''; kepada pembaca kelas lima pada bab
\omterm{def:g5:everyday-energy:energy}{energinya} --- ``rantai yang kamu telusuri akan menjadi rumus yang
berisi \omterm{def:g9:kinetic-energy-safety:joule}{joule}''; kepada pembaca kelas delapan pada bab hukum Ohmnya ---
``potret yang kamu gambar adalah kebiasaan menguji hukum, dan itulah
seluruh rahasianya''.)
\end{solution}

\begin{solution}{pb:g9:atoms-to-galaxies:1}
\textbf{1.} Halaman sekolahnya: sekon $2$; Bumi: sekon $7$; \omterm{def:g6:sun-earth-moon:axisorbit}{orbit}
\omterm{def:g2:moon-in-the-sky:moon}{Bulan}: sekon $9$.
\textbf{2.} Sekon $11$ (pada $1.5 \times 10^{11}$ \unit{m});
penceritanya: ``sinar matahari di atas tikarnya berumur delapan
menit.''
\textbf{3.} Kasarnya sekon $13$ sampai $16$: bingkainya memuat
seluruh \omterm{def:g4:solar-system:family}{Tata Suryanya} sebagai sebuah titik dan belum ada \omterm{def:g4:solar-system:star}{bintang} ---
dengan jujur, kekosongan yang hampir hitam, beberapa sekon lamanya:
babak filmnya yang paling benar.
\textbf{4.} Langit tersurvei yang terdalam, $10^{26}$ \unit{m} berisi
galaksi --- dengan \omterm{def:g1:light-and-shadows:source}{cahaya} tertuanya lebih tua daripada Bumi sendiri.
\textbf{5.} Semutnya: sekon $2$ ke dalam; selnya: sekon $5$;
\omterm{def:g7:states-of-matter:molecule}{molekulnya}: sekon $9$; \omterm{def:g9:atoms-to-galaxies:atom}{atomnya}: sekon $10$.
\textbf{6.} Menembus kekosongan milik \omterm{def:g9:atoms-to-galaxies:atom}{atomnya} sendiri --- lima sekon
ketiadaan di antara awannya dan pusatnya: ``kalau \omterm{def:g9:atoms-to-galaxies:atom}{atom} ini sebuah
stadion, kita sudah meninggalkan tribun yang tertinggi dan masih
terbang menuju sebutir kacang polong.''
\textbf{7.} \omterm{def:g9:atoms-to-galaxies:atom}{Intinya} --- dan dengan jujur: ``di sini tangga kami
berakhir, bukan karena alamnya berhenti, melainkan karena inilah anak
tangga terkecil yang pernah diukur manusia; bangunannya berjalan
terus.''
\textbf{8.} Ke luar $26$ sekon; ke dalam $15$: sebuah alam semesta
empat puluh satu sekon.
\textbf{9.} Tiap sekonnya mengalikan pemandangannya sepuluh kali lipat
--- \emph{nisbah} yang sama, bukan penambahan yang sama: logika
tangganya adalah bahwa satu anak tangga berupa sebuah faktor, dan
menurut faktornya langkah jalanan dan langkah gugusnya sama besar.
\textbf{10.} Misalnya: ``Sepenuh apa pun tampaknya, hampir segala yang
ada di sini kosong --- materi dan langit sama saja berupa pulau yang
langka di dalam ketiadaan yang luas.''
\textbf{11.} Mikroskopnya (bab \omterm{def:g8:lenses-images:families}{lensanya}: dua \omterm{def:g8:lenses-images:families}{lensa cembung} di atas
yang kecil) dan teleskopnya (bab yang sama, yang mengumpulkan yang
samar dan jauh) --- kunci kedua arahnya, yang keduanya diasah dari
optika buku ini.
\textbf{12.} Misalnya: ``Sembilan tahun, satu tangga: dari kaki seekor
kepik yang dapat dihitung sampai galaksi yang lebih tua daripada tanah
di bawah kaki, setiap anak tangganya dipasang oleh pengukuran yang
boleh diperiksa siapa pun. Tahun depan pertanyaan yang kami tinggalkan
terbuka --- \omterm{def:g2:pushes-and-pulls:force}{gaya} dan gerak, hakikat \omterm{def:g1:light-and-shadows:source}{cahaya}, jantung \omterm{def:g9:atoms-to-galaxies:atom}{atomnya} --- mulai
terjawab; bawalah tangganya.''
\end{solution}
